Every quantum algorithm for a differential equation pays a ratio that its discretisation fixes before any circuit is built: the condition number of a linear system, the norm of a generator times the evolution time, the ratio of nonlinearity to dissipation in a lift, or the norm of an encoded Hamiltonian against its gap. An accurate discretisation can still make that ratio large, and it typically grows with the resolution. The review follows the ratio through three bodies of work: the quantum routes to differential equations, judged by what each assumes, which ratio it pays, what it returns and where it fails; the classical methods the thesis adapts, the spectral discretisations on which the physics-informed construction rests, for which numerical analysis has already shown how the choice of basis and of unknowns sets the ratio and how it can be lowered, and Carleman's linearisation of a nonlinear equation; and the quantum primitives that turn a ratio into a count of queries, the filter degree of \Cref{eq:degree}. The \gls{pihm} itself is presented and judged separately, in \Cref{sec:pihm}.

\subsection{Quantum Computation Preliminaries}
\label{sec:bg-qc}

We use the circuit model of Nielsen and Chuang \cite{nielsenchuang2010}. The state of $n$ qubits is a unit vector in $\mathbb{C}^N$, $N=2^n$, and a vector $v\in\mathbb{C}^N$ is \emph{amplitude-encoded} as $\ket{v}=v/\|v\|$. The encoding stores only the direction of $v$, so its length must be carried separately (for a differential equation, the scale $\eta$ of \Cref{sec:pihm-model}), and it gives access to $v$ only through measurement, which returns the index $j$ with probability $|v_j|^2/\|v\|^2$; recovering all $N$ entries takes at least of order $N$ repetitions (\Cref{sec:bg-filtering}). Circuits of unitary gates act on these states, and a non-unitary operation is applied by acting on the system together with \emph{ancilla} qubits prepared in $\ket{0}$ and keeping only the outcomes in which they are measured in $\ket{0}$ again. This \emph{post-selection} succeeds with a probability that enters the cost; every non-unitary matrix of the thesis, from a residual to a filter, acts on a state this way (\Cref{sec:bg-block-encoding}).

An algorithm that accesses a matrix only through a circuit encoding it is costed in \emph{queries}, the uses of that circuit, and its gate count is the number of queries times the gates per query, on a register of system and ancilla qubits. The asymptotic results reviewed below, and the filter degree of \Cref{eq:degree}, count queries; the gates per query belong to the encoding circuit. \Cref{sec:standard} separates the two, and \Cref{sec:pipeline-design} fixes how gates are counted. A tilde on $\mathcal O$, $\Omega$ or $\Theta$ suppresses factors logarithmic in $N$ and in $1/\varepsilon$, and a cost polynomial in $n=\log_2N$ is called poly-logarithmic.

Every cost in the thesis is that of a fault-tolerant machine, whose logical qubits and gates are protected by error correction. Current \gls{nisq} processors have no error correction, and noise limits the depth of a reliable circuit \cite{preskill2018nisq}; the variational algorithms of \Cref{sec:bg-de-variational} target them. A filter applies its block encoding as many times in sequence as its degree, which grows with the resolution, so its feasibility is judged against fault-tolerant resource estimates (\Cref{sec:disc-feasibility}) and NISQ hardware is out of scope.

\subsection{Quantum Algorithms for Differential Equations}
\label{sec:bg-de}

Aaronson's caveats on the linear-systems algorithm apply to every route below: the input must load efficiently, the operator must be well conditioned and efficiently accessible, and the output is a quantum state, useful only where a few properties of the solution suffice \cite{aaronson2015fineprint}. \Cref{tab:de-approaches} compares the routes under these headings.

\begin{table}[tbp]
    \centering
    \caption[Quantum routes to differential equations and what sets their cost]{\textbf{Each quantum route's cost is set by a ratio its discretisation fixes.} Rows: the routes of \Cref{sec:bg-de}; columns: assumptions, cost driver, output, failure mode. $A$, the generator of a linear system; $\kappa$, its condition number; $T$, evolution time; $x(t)$, the solution at time $t$; $\alpha_H$, $\Delta$, subnormalisation and spectral gap of a block-encoded Hamiltonian $H$.}
    \label{tab:de-approaches}
    \footnotesize
    \begin{tabular}{@{}p{0.16\linewidth}p{0.20\linewidth}p{0.19\linewidth}p{0.14\linewidth}p{0.21\linewidth}@{}}
        \toprule
        Route & Assumes & Cost depends on & Returns & Fails where \\
        \midrule
        Linear systems \cite{harrow2009hhl,berry2017linearode,childsliu2020spectral,costa2022optimal,jennings2023cost} & linear equation in first-order form; efficient access to $A$ and loading of the data & $\kappa$ of the discretised system, which grows with $T$ & the solution, or its history, as a state & $\kappa$ grows with the resolution; readout; at most polynomial gain at fixed dimension \\
        \addlinespace
        Hamiltonian simulation; linear combination of Hamiltonian simulations; Schrödingerisation \cite{lowchuang2019qubitization,anliulin2023lchs,jinliuyu2023schrodingerisation} & linear dynamics; for the linear combination, a positive semi-definite dissipative part & $\|A\|T$; the decay $\|x(0)\|/\|x(T)\|$ & the state at time $T$ & nonlinear dynamics; $\|A\|$ grows with the resolution; readout \\
        \addlinespace
        Variational \cite{lubasch2020variational,kyriienko2021differentiable,bravoprieto2023vqls,pool2024nonlinear} & a trainable parametrised circuit & optimisation steps times samples, unbounded & a trained circuit & no guarantee of convergence; barren plateaus \\
        \addlinespace
        Ground-state preparation \cite{subasi2019adiabatic,anlin2022timeoptimal,costa2022optimal,lintong2020filtering} & a positive semi-definite $H$ whose null vector is the solution; a known gap; an initial state of non-vanishing overlap & $\alpha_H/\Delta$, or its square root for a filter at the spectral edge; for a linear system, $\Delta$ of order $1/\kappa$ & the solution as a ground state & the gap closes as the problem's conditioning grows; readout \\
        \bottomrule
    \end{tabular}
\end{table}

\paragraph{Linear systems.}
\label{sec:bg-de-linear}

The first route discretises the equation, in time and where needed in space, into one linear system whose solution holds the trajectory, and solves that system with a quantum linear-systems algorithm. The \gls{hhl} algorithm prepares a state proportional to $A^{-1}b$ for an $s$-sparse $N\times N$ matrix of condition number $\kappa$ in time $\tilde{\mathcal O}(s^2\kappa^2\log N/\varepsilon)$, and shows that time $\kappa^{1-\delta}$ for any $\delta>0$ would imply BQP $=$ PSPACE \cite{harrow2009hhl}. Discretising a differential equation into such a system gave the first solvers: Berry with a linear multistep method, at a cost polynomial in $1/\varepsilon$ \cite{berry2014highorder}; Berry, Childs, Ostrander and Wang with a truncated Taylor series, at a cost polynomial in $\log(1/\varepsilon)$ \cite{berry2017linearode}; and Childs and Liu with Chebyshev collocation in time, whose collocation matrix is dense and whose condition number they bound by a polynomial of degree $3.5$ in the number of nodes \cite{childsliu2020spectral}. Childs, Liu and Ostrander extended finite-difference and spectral discretisations to \glspl{pde} at the same dependence on precision \cite{childsliuostrander2021pde}, Costa et al.\ brought the linear-systems solver to the optimal $\mathcal O(\kappa\log(1/\varepsilon))$ with a discrete adiabatic theorem \cite{costa2022optimal}, and Jennings et al.\ gave the first non-asymptotic costs, constant factors included, preparing the history state of any stable linear system with complexity $\mathcal O(T^{1/2})$ in the evolution time, a fast-forwarding that classical stability permits \cite{jennings2023cost}. The ordinary-equation solvers take the equation as a first-order system $x'=Ax+b$, reducing any higher order before discretising, so the conditioning they pay is that of a first-order operator. The route's strength is its generality at a proven cost: any linear equation whose discretised operator can be accessed efficiently is solved at a cost logarithmic in $N$ and, from Berry et al.\ on, in $1/\varepsilon$. Its weakness is that the cost runs through $\kappa$, which the discretisation fixes and which grows with the resolution and with the evolution time, so the discretisation is part of the cost (\Cref{sec:bg-spectral}). Learning all $N$ amplitudes of the output takes at least of order $N$ repetitions, and for equations in a fixed number of dimensions with smooth solutions the advantage over the best classical methods is at most polynomial \cite{montanaro2016fem}.

\paragraph{Hamiltonian simulation and Schrödingerisation.}
\label{sec:bg-de-hamsim}

A linear equation $x'=-\mathrm{i}Hx$ with Hermitian $H$ is a Schrödinger equation, and integrating it to time $T$ is Hamiltonian simulation, which qubitisation performs with $\mathcal O(\alpha T+\log(1/\varepsilon))$ queries to a block encoding of $H$ at subnormalisation $\alpha$, an optimal dependence on both \cite{lowchuang2017qsp,lowchuang2019qubitization}. Two constructions carry non-unitary linear dynamics $x'=-Ax$ to this primitive. The linear combination of Hamiltonian simulations writes the propagator of a dissipative system as a weighted integral of unitary evolutions and implements it as a \gls{lcu}, without dilating the problem into a larger linear system and at optimal state-preparation cost \cite{anliulin2023lchs}. Schrödingerisation maps a linear ordinary or partial differential equation to a Schrödinger equation in one more dimension by a warped phase transformation, and recovers the solution from the auxiliary variable \cite{jinliuyu2023schrodingerisation}. Neither solves a linear system, so neither pays a condition number: the cost runs through $\|A\|T$, the norm of the generator times the evolution time. The route is linear, and its ratio also grows with the resolution, since the norm of a discretised spatial operator grows with the number of unknowns, as the square of the resolution for a second derivative. It returns the state at the final time only, and recovering it from a dissipative evolution needs a number of repetitions that grows with the decay $\|x(0)\|/\|x(T)\|$ of the solution's norm. The linear combination, as stated, also requires the dissipative part of $A$, its Hermitian part, to be positive semi-definite, which restricts it to stable dynamics.

\paragraph{Variational algorithms.}
\label{sec:bg-de-variational}

A variational algorithm represents the solution by a parametrised circuit and minimises a cost built from the equation with a classical optimiser. Lubasch et al.\ evaluate a nonlinear cost on several copies of the state, as demonstrated on the nonlinear Schrödinger equation \cite{lubasch2020variational}. The differentiable quantum circuits of Kyriienko, Paine and Elfving represent the solution as the expectation value of a circuit with input $x$, differentiate it with respect to $x$ by parameter shifts on the feature map, and train it on the residual of the equation and its data \cite{kyriienko2021differentiable}, the loss of a \gls{pinn} \cite{raissi2019pinn}. The variational linear solver finds a ground state the same way, minimising $\braket{x|A^\dagger(I-\ket{b}\bra{b})A|x}$, whose zero-energy ground state solves $Ax\propto b$ \cite{bravoprieto2023vqls}; encoding time in a register, as the Feynman--Kitaev clock does, puts a whole trajectory in one ground state \cite{barison2022variational}, and a space--time encoding of this kind has given the evolution of Burgers' equation from a single ground-state computation \cite{pool2024nonlinear}. The route's strengths are its circuits, shallow enough for NISQ devices, and its reach: a nonlinear term enters the cost directly, with no linearisation. Its weakness is that nothing bounds the number of optimisation steps or guarantees the optimum. The gradients of a sufficiently random parametrised circuit vanish exponentially in the number of qubits, the barren plateau \cite{mcclean2018barren}, and each evaluation of the cost is a sampled expectation value, so the cost of a solution is unknown before it is found.

\paragraph{Ground-state preparation.}
\label{sec:bg-de-ground}

The last route builds a positive semi-definite Hamiltonian whose null vector is the solution and prepares that vector without optimisation. Subaşı, Somma and Orsucci do so for a linear system by randomised adiabatic evolution at cost $\mathcal O(\kappa\log\kappa/\varepsilon)$ \cite{subasi2019adiabatic}, An and Lin by a time-optimal adiabatic schedule \cite{anlin2022timeoptimal}, and Costa et al.\ by a discrete adiabatic walk \cite{costa2022optimal}. In Lin and Tong's form \cite{lintong2020filtering}, for positive-definite $A$ with $\|A\|=1$ and $Q_b=I-\ket{b}\bra{b}$, the Hamiltonian with off-diagonal blocks $AQ_b$ and $Q_bA$ has $\ket{0}\ket{x}$ in its null space beside $\ket{1}\ket{b}$, which a measurement of the first qubit removes; the rest of its spectrum lies at a distance of order $1/\kappa$ from zero, and an optimal Chebyshev filter projects onto the null space. The route's strength is a proven cost with nothing to train, in a form that separates the circuit from the spectrum: a filter's degree depends on the Hamiltonian only through the subnormalisation of its block encoding and its gap, and its error falls exponentially with the degree (\Cref{sec:bg-filtering}). Its weakness is that the gap carries the problem's conditioning, of order $1/\kappa$ for a linear system, so the route pays the ratio of the linear-systems route in another form. A filter also needs the gap, or a bound on it, in advance, and an initial state whose overlap with the solution does not vanish. The physics-informed effective Hamiltonian of Wu et al.\ belongs to this route: it builds the Hamiltonian from a Chebyshev discretisation of the equation and its data and prepares the ground state with a filter (\Cref{sec:pihm}).

\paragraph{Critical assessment.}
\label{sec:bg-de-assessment}

Each route in \Cref{tab:de-approaches} pays a ratio set by how the equation is discretised and encoded, which an accurate discretisation can still make large: the condition number of a linear system, the norm of a generator times the evolution time, or the norm of an encoded Hamiltonian against its gap. Every route returns a quantum state, so its advantage holds only where a few properties of the solution suffice \cite{aaronson2015fineprint,williams2024readout}, and is at most polynomial in a fixed, low number of dimensions \cite{montanaro2016fem}. The thesis therefore compares quantum methods with one another, at equal accuracy, in queries and gates, and claims no advantage over classical solvers. The guarantees differ: the linear-systems, Hamiltonian-simulation and ground-state routes have proven costs, and the variational route has none. Only the variational route takes a nonlinear equation as it stands; the others are linear and reach a nonlinear equation through a linearisation such as Carleman's, which adds a ratio of its own (\Cref{sec:bg-carleman}).

\subsection{Classical Discretisation and Linearisation}
\label{sec:bg-spectral}

The ratios of \Cref{tab:de-approaches} are set when the equation is discretised. For the Chebyshev discretisation on which the physics-informed construction is built, classical numerical analysis locates the source of the ratio, the density of differentiation in coefficient space, and two ways of lowering it: a basis in which differentiation is banded, and derivatives carried as unknowns. A nonlinear equation reaches the linear routes only after a second classical step, Carleman's linearisation, which brings a ratio of its own.

\paragraph{Chebyshev discretisation and differentiation.}
\label{sec:bg-chebyshev}

A function analytic on $[-1,1]$ has a Chebyshev series $f=\sum_kb_kT_k$ with geometrically decaying coefficients, so the $N$-term truncation converges faster than any power of $N$ \cite{trefethen2013approximation}; the physics-informed construction, like the spectral methods it draws on, works with these coefficients for that reason. A differential equation becomes a finite system by Galerkin projection onto the basis, by the tau method, which projects but replaces the highest-degree rows by the boundary conditions, or by collocation at a set of nodes \cite{canuto2006spectral,boyd2001Chebyshev}. Wu et al.'s construction is closest to the tau method, with the equation acting on the coefficients and the data imposed as rows at points, but its system is overdetermined and solved by least squares (\Cref{sec:pihm-residual}).

Differentiation in coefficient space is exact on polynomials of degree below $N$ and dense, because the derivative of $T_p$ contains every $T_k$ of the opposite parity below it: the matrix is upper triangular, about a quarter of its entries are non-zero, and its norm grows as $N^2$ (\Cref{sec:standard-operators}). An operator of order $k$ built from its powers grows as $N^{2k}$, and so, typically, does the condition number of a tau or collocation discretisation \cite{olver2013,canuto2006spectral}.

\paragraph{The ultraspherical method.}
\label{sec:bg-ultraspherical}

Olver and Townsend make differentiation banded by letting the derivative live in another basis \cite{olver2013}. Since $T_k'=k\,C^{(1)}_{k-1}$, with $C^{(1)}_k=U_k$ the Chebyshev polynomials of the second kind, differentiation from Chebyshev to $C^{(1)}$ coefficients is a single superdiagonal, and conversion of a series' own coefficients to $C^{(1)}$ has two non-zero diagonals. Raising the basis one ultraspherical (Gegenbauer) level $C^{(\lambda)}$ per derivative gives, for an equation of order $k$ with polynomial coefficients, a leading coefficient that does not vanish and $k$ boundary conditions, an operator banded apart from its boundary rows, whose condition number grows linearly with the truncation and is bounded after a diagonal preconditioner, in the 2-norm for Dirichlet conditions \cite{olver2013}. The dense differentiation matrix we use factorises as $\G=\Bt^{-1}\Dt$, with $\Bt$ and $\Dt$ the method's first conversion and differentiation operators up to a factor of two (\Cref{app:ultraspherical}); the descriptor form of \Cref{sec:descriptor} keeps the two factors, and its block rescaling takes the place of the preconditioner.

Olver and Townsend eliminate the lower derivatives and solve one banded system for the solution's coefficients, whereas the descriptor keeps every derivative as an unknown in the Chebyshev basis, tied to the next by the same pair $(\Bt,\Dt)$ (\Cref{sec:bg-fosls}). The full reduction, banded on $N$ unknowns without a block register, is not compared here, although its conditioning suggests, untested, that it would match the descriptor's degree on a regular problem (\Cref{sec:limitations}).

\paragraph{First-order system least squares.}
\label{sec:bg-fosls}

The normal equations of a least-squares method square the residual operator, so high order compounds the conditioning: least squares applied directly to a second-order elliptic equation needs continuously differentiable finite elements and gives condition numbers of order $h^{-4}$ in the mesh width $h$, against $h^{-2}$ for a Galerkin method \cite{bochevgunzburger1998review}. \Gls{fosls} introduces derivatives as unknowns, such as a flux $v=\nabla\phi$, and minimises the residuals of the resulting first-order system \cite{cai1994fosls}. Where the functional is equivalent to a product of $H^1$ norms, the condition numbers return to order $h^{-2}$ and the unknowns nearly decouple, so block-diagonal preconditioners are effective \cite{bochevgunzburger1998review,bochev2009lsfem}.

The physics-informed Hamiltonian is such a functional. Stacking the operators of the equation and of its data into one residual operator $R$ gives $H=R^\dagger R$, whose energy $\|R\psi\|^2$ is the squared residual of the equation and its data (\Cref{sec:pihm-residual,sec:standard-form}). Eliminating every derivative through $\G$, as Wu et al.\ do, is the high-order least-squares method; carrying the derivatives as unknowns is FOSLS in a spectral basis. The quantum ordinary-equation solvers of \Cref{sec:bg-de-linear} already work with first-order systems, and the physics-informed construction departs from them in eliminating the derivatives into one dense operator. A filter applied through a block encoding of $R$ pays the conditioning of $R$ (\Cref{sec:bg-filtering}); the squaring falls only on an eigensolver acting on $H$. A first-order formulation still lowers what the filter pays, because each block of $R$ then carries one derivative instead of $k$ (\Cref{sec:standard-hamiltonian}).

\paragraph{What we add.} Both ideas are classical. The increment is their transfer to a block-encoded physics-informed Hamiltonian and its consequences for preparing the ground state: a proof that the reformulated residual has the standard form's kernel, so that the prepared state is the same solution (\Cref{sec:descriptor-layouts}); an encoding of the banded residual as one flat linear combination of unitaries with $\mathcal O(\log n)$ ancilla qubits (\Cref{sec:descriptor-block-encoding}); the finding that the carried derivatives erode the Hamiltonian's gap, and a block rescaling that restores it on one axis without ancilla and not on two (\Cref{sec:descriptor-gap,sec:descriptor-rescaling}); and a measurement of what the reformulation does to the filter's degree, as built and at an ideal encoding (\Cref{sec:results-scaling,sec:results-gap}).

\paragraph{Fourier--Galerkin truncation.}
\label{sec:bg-fourier}

On a periodic domain the Fourier basis diagonalises differentiation, $\partial_xe^{\mathrm{i}kx}=\mathrm{i}k\,e^{\mathrm{i}kx}$, so linear terms act mode by mode and a quadratic term becomes a convolution of coefficients \cite{canuto2006spectral}. A Galerkin truncation keeps exactly the convolution terms whose output lies in the retained set of modes; a pseudo-spectral evaluation on a grid aliases the discarded terms onto retained modes unless it is dealiased \cite{canuto2006spectral}. The fluids of \Cref{sec:fluids} use Galerkin truncations on real Fourier modes, so their quadratic coupling is a fixed matrix of exact coefficients, as the Carleman lift requires; a pseudo-spectral simulation serves only as the reference for the truncation's error.

\paragraph{Carleman linearisation.}
\label{sec:bg-carleman}

Carleman's embedding replaces a polynomial system of ordinary differential equations by an infinite linear system for the monomials of its unknowns \cite{carleman1932,kowalskisteeb1991}. For $u'=F_1u+F_2u^{\otimes2}$ the vector $(u,u^{\otimes2},u^{\otimes3},\dots)$ obeys a linear equation in which each level couples only to itself and to the next, and truncation at order $K$ leaves a finite linear system whose first level approximates $u$. Its use in quantum algorithms followed an earlier approach of Leyton and Osborne, which transforms the amplitudes of many copies of the state nonlinearly, at resources exponential in the integration time \cite{leyton2008nonlinear}. Liu et al.\ solve the truncated Carleman system with a quantum linear-systems algorithm for dissipative quadratic equations in $M$ variables, at cost $T^2q\,\mathrm{poly}(\log T,\log M,\log1/\varepsilon)/\varepsilon$, with $q$ the decay of the solution, when the ratio $R$ of nonlinearity and forcing to dissipation satisfies $R<1$, and prove the problem intractable in the worst case for $R\ge\sqrt2$ \cite{liu2021carleman}; Krovi makes the dependence on precision logarithmic and relaxes the normality of the linear part to a negative log-norm \cite{krovi2023improved}.

The lift is linear and exact before truncation, so a nonlinear equation inherits the guarantees of whichever linear route solves the truncated system, and the dimension $M^K$ of the lift fits in about $K\log_2M$ qubits. Its weakness is the truncation. The condition $R<1$ bounds the truncation error uniformly over states, so it is sufficient and not necessary: a lift may converge for $R>1$, which \Cref{sec:carleman} (writing the ratio $\rho$) measures and \Cref{sec:results-nonlinear} reports. Dissipation must still dominate, since the truncation is controlled only where the nonlinearity is weak against the damping, and the lift's dimension, of order $M^K$, enters the cost of every route that solves it.

\subsection{Block Encodings, Signal Processing and Filtering}
\label{sec:bg-primitives}

A block encoding applies a non-unitary matrix on a quantum computer, a polynomial transformation of the block encoding builds a filter, and the filter's degree counts the queries that a prepared ground state costs.

\paragraph{Block encodings and linear combinations of unitaries.}
\label{sec:bg-block-encoding}

\begin{definition}[Block encoding \cite{gilyen2019qsvt}]
\label{def:block-encoding}
A unitary $U$ on $m+n$ qubits is an $(\alpha,m,\epsilon)$ \emph{block encoding} of an $N\times N$ matrix $A$ if
\begin{equation}
    \bigl\|A-\alpha\,(\bra{0^m}\otimes I)\,U\,(\ket{0^m}\otimes I)\bigr\|\le\epsilon ;
    \label{eq:block-encoding}
\end{equation}
$\alpha$ is its \emph{subnormalisation}, and the encoding is \emph{exact} when $\epsilon=0$.
\end{definition}

Applied to $\ket{0^m}\ket{\psi}$ and post-selected on the ancillas' zero, an exact block encoding applies $A/\alpha$ with success probability $\|A\psi\|^2/\alpha^2$. The block of a unitary has norm at most one, so $\alpha\ge\|A\|$ for every exact block encoding, and the excess $\alpha/\|A\|$ multiplies the cost of every algorithm that uses it. The standard construction is the \gls{lcu} \cite{childs2012lcu}: for $A=\sum_{t<L}c_tU_t$ with unitaries $U_t$, PREPARE loads $\sum_t\sqrt{|c_t|/\alpha}\,\ket{t}$ on $\lceil\log_2L\rceil$ ancilla qubits, SELECT applies $U_t$ controlled on $\ket{t}$, with the phases of the $c_t$ absorbed into the $U_t$, and PREPARE$^\dagger$ returns the ancillas, giving an exact block encoding at $\alpha=\sum_t|c_t|$. An exact combination whose SELECT indexes $L$ distinct operations needs $\lceil\log_2L\rceil$ selector qubits \cite{chakraborty2025qsvt}. Explicit circuits exist for many structured and sparse matrices \cite{camps2024explicit,sunderhauf2024blockencoding}; each carries its own $\alpha$, and the two forms of the residual compared in \Cref{sec:results} differ in both $\|A\|$ and $\alpha/\|A\|$.

If $U$ is Hermitian and block-encodes $H/\alpha$, the walk $W=(2\Pi-I)\,U$, with $\Pi=\ket{0^m}\bra{0^m}\otimes I$, satisfies $\Pi W^k\Pi=T_k(H/\alpha)$: its $k$-th power block-encodes the $k$-th Chebyshev polynomial of the encoded matrix at one query per power \cite{lowchuang2019qubitization}, so a polynomial of degree $d$ in $H/\alpha$ costs $d$ queries once its combination is formed.

\paragraph{Quantum signal processing.}
\label{sec:bg-qsp}

Quantum signal processing interleaves a single-qubit signal rotation with fixed phase rotations to realise a polynomial of the signal \cite{lowchuang2017qsp}. The \gls{qsvt} lifts the construction to block encodings, applying a degree-$d$ polynomial $P$ to the singular values of the encoded matrix with $d$ queries and one further ancilla \cite{gilyen2019qsvt}, and places search, Hamiltonian simulation and linear systems in one framework \cite{martyn2021grand}. Its polynomials are bounded by one on $[-1,1]$ and of definite parity, so a function of indefinite parity, such as the imaginary-time filter $e^{-\beta x}$, is split into even and odd parts, each with its own phase sequence, and the two are summed by a further linear combination at a cost in subnormalisation; Wu et al.\ use this mixed-parity construction (\Cref{sec:pihm-encoding}). \Gls{gqsp} drops the parity constraint \cite{motlagh2024gqsp}. SU(2) rotations on one signal qubit, with the encoded unitary applied controlled on it, realise any $P(z)=\sum_{k\le d}q_kz^k$ with $|P|\le1$ on the unit circle at $d$ queries, and applied to the walk $W$ they give $\sum_kq_kT_k$ of the encoded matrix (\Cref{sec:pipeline-filter}). The phases follow from a complementary polynomial, computed by a fast-Fourier-transform formula with explicit error guarantees \cite{berntson2025complementary}, through the inverse of the \gls{nlft} to which GQSP is equivalent \cite{laneve2025gqspnlft,nlfft2025su2}.

\paragraph{Ground-state preparation by filtering.}
\label{sec:bg-filtering}

A filter is a function of $H$ that keeps the ground eigencomponent of an initial state $\psi_{\mathrm{init}}$ and suppresses the rest. Imaginary-time evolution $e^{-\beta H}$ does so once $\beta$ exceeds the inverse gap \cite{motta2020qite}, and inverse iteration $H^{-k}$ by powers \cite{kyriienko2020inverse}; expanded in polynomials, both are implementable by the transformations above. For an eigenvalue at zero with the rest of the spectrum in $[-1,-\Delta]\cup[\Delta,1]$, Lin and Tong's eigenstate filter, a Chebyshev polynomial in $x^2$, has the least maximum on that set among polynomials of its degree equal to one at zero, and needs degree $\mathcal O(\Delta^{-1}\log(1/\varepsilon))$ \cite{lintong2020filtering}. Their ground-state algorithm approximates a shifted sign function, given a bound that separates the ground energy from the rest of the spectrum, with $\mathcal O\bigl((\alpha/(\gamma\Delta))\log(1/\varepsilon)\bigr)$ queries for subnormalisation $\alpha$, gap $\Delta$ and overlap $\gamma=|\braket{\psi_0|\psi_{\mathrm{init}}}|$ with the ground state; they argue, by reduction from unstructured search, that its dependence on $\gamma$ and $\Delta$ is essentially optimal \cite{lintong2020groundstate}. Both assume the gap, or bounds that imply it, known in advance, a standard assumption that we share (\Cref{sec:pipeline-filter}).

Both place the ground eigenvalue inside the encoded spectrum; at its edge the degree falls quadratically. Let $H\succeq0$ be block-encoded with its spectrum mapped into $[-1,1]$ and the ground value at the lower edge, as by $2H/\alpha_H-I$, with gap $\Delta=\lambda_1-\lambda_0$ and a ground value small against $\Delta$. The excited spectrum lies in $[x_1,1]$, at a distance of order $\Delta/\alpha_H$ from the ground value. Among polynomials of degree $d$ bounded by one on $[-1,1]$, the Chebyshev polynomial of $[x_1,1]$ has the least maximum there relative to its value at the ground state, by Chebyshev's classical inequality \cite{borweinerdelyi1995}; it is the edge counterpart of Lin and Tong's filter. Outside its interval it grows as the hyperbolic cosine of the degree times the square root of the distance, so suppressing the excited spectrum to $\varepsilon$ relative to the ground state needs the degree
\begin{equation}
    d = \Theta\!\left(\sqrt{\frac{\alpha_H}{\Delta}}\,\log\frac{1}{\varepsilon}\right),
    \qquad
    \alpha_H = \alpha_R^2 .
    \label{eq:degree}
\end{equation}
The square root comes from the ground value's place at the edge of the encoded spectrum, where a Chebyshev polynomial grows as the square root of the distance; for frustration-free Hamiltonians, sums of positive terms that all vanish on the ground state, it is Somma and Boixo's spectral gap amplification \cite{somma2013gap}, and $H=R^\dagger R$ is one exactly when its kernel is exact. A small non-zero ground value leaves the degree unchanged, since the degree is set by the gap, and costs only the square of the filter's value at the ground state, a factor close to one, in the probability of success (\Cref{sec:pipeline-filter}, \Cref{app:minimax-filter}). Here $\alpha_R$ is the subnormalisation of a block encoding of $R$, from which one reflection block-encodes $H$ at $\alpha_H=\alpha_R^2$ (\Cref{sec:pipeline-walk}); the walk's eigenphases are set by the singular values of $R/\alpha_R$, and the gap the filter resolves is of order $\sigma_1/\alpha_R=(\Delta/\alpha_H)^{1/2}$ (\Cref{sec:pipeline-walk}).

Each degree is one query, so $d$ counts the queries of one filtered preparation before amplification. The circuit enters only through $\alpha_H\ge\|H\|$ (\Cref{sec:bg-block-encoding}), so with this filter no encoding brings the degree below order $\sqrt{\|H\|/\Delta}\,\log(1/\varepsilon)$. With $\alpha_H=\alpha_R^2$ the degree is linear in $\alpha_R$; for an exact kernel, whose gap is the square of the smallest non-zero singular value of $R$, $\sqrt{\alpha_H/\Delta}$ is at least the ratio of the largest singular value of $R$ to that one, the residual's effective condition number. The cost of a filtered ground state is therefore set by the norm of the operator it is built from, measured against its gap. The bound holds for polynomial filters given the edge and the gap, among which Chebyshev's inequality makes the degree least; over all algorithms the thesis claims no more than Lin and Tong's optimality supports (\Cref{sec:disc-complexity}).

Post-selection succeeds with probability of order $\gamma^2$, and amplitude amplification raises it to near one with $\mathcal O(1/\gamma)$ repetitions \cite{brassard2000amplitude}, so a preparation costs of order $d/\gamma$ queries; an answer-independent initial state with non-vanishing overlap is part of the method (\Cref{sec:pipeline-filter}).

\paragraph{Readout.} An overlap or expectation value is estimated to additive error $\delta$ by a Hadamard or swap test with $\mathcal O(1/\delta^2)$ repetitions, or with $\mathcal O(1/\delta)$ by amplitude estimation, including variants without phase estimation \cite{brassard2000amplitude,suzuki2020amplitude,grinko2021iterative}; a full field needs of order $N$ such estimates, the readout problem of quantum differential-equation algorithms \cite{williams2024readout,nishi2025overlap}. Wu et al.\ read the field point by point with an interferometric overlap measurement (\Cref{sec:pihm-encoding}). We decode the field classically from the prepared state, which our classical computations make available (\Cref{sec:pipeline-readout}), and leave readout unadressed, following the convention of every route of \Cref{tab:de-approaches}.

\subsection{Synthesis}
\label{sec:bg-synthesis}

The physics-informed effective Hamiltonian sits where the three strands meet. It is a ground-state route, so it inherits the filters' guarantees and their cost law, \Cref{eq:degree}: the square root of an encoded Hamiltonian's norm over its gap. It is a least-squares functional on a Chebyshev discretisation, so it inherits spectral accuracy together with the conditioning of dense differentiation, the property that the ultraspherical method and first-order-system least squares each address. And it descends from physics-informed learning but trains nothing. Its authors name the gap's behaviour and the cost of differentiation as questions for future work; in the terms of \Cref{eq:degree} the first is the degree's denominator and, as \Cref{sec:pihm-assessment} argues, the second sets its numerator. \Cref{sec:pihm} presents the method on the same attributed footing as the material above, judges it, and poses those questions in the form we answer.
